\subsection{The finite CM calculation and its independent verification}

In this section $B=\langle s,r\mid s^2=r^3=1\rangle$ is a discrete
free product, $t=sr$, and $F=\ker(B\to S_3)$ is the discrete free
group on $x=u^2$, $y=v^2$, where $u=t^{-1}$ and $v=sus$.
Let $q:B\twoheadrightarrow Q=\mathrm{SL}_2(\Z/12)/\{\pm I\}$
send $s,r$ to the classes of
\[
 S=\begin{pmatrix}0&1\\-1&0\end{pmatrix},\qquad
 R=\begin{pmatrix}0&-1\\1&-1\end{pmatrix}.
\]
Write $N_B=\ker q$, $C=N_B^{\mathrm{ab}}$, and
$K_j=[N_B,N_B]N_B^j$, so $G_j=B/K_j$. All these maps are homomorphisms
of discrete groups; their maps to finite targets extend continuously
to the profinite completions. The inverse image of $K_j$ under
$B_3\twoheadrightarrow B$ is $\widetilde N_j$.
The CM comparison is LM7 D.15--D.19 and D.25--D.26, in the setting
of \emph{La Longue Marche} \S51.15, scans 156--170. The assertions
below concern these specified finite coordinates.

\begin{statement}{E.1}{Proved}{Independent confirmation of Task E.1}
The order of $Q$ is $576$, and $C\simeq\Z^{97}$. With the CM sign
convention of LM7 D.25, the independent calculation gives
\[
\begin{array}{l|rr}
 &\mathrm{CM\ compatible}&\mathrm{extra\ CM\ coset}\\\hline
 \mathrm{matrix\ pairs}&24&24\\
 \mathrm{cusp\ equation\ in}\ G_2&12&12\\
 \mathrm{cusp\ equation\ in}\ G_3&24&24\\
 \mathrm{certified\ automorphisms\ of}\ G_6&12&12.
\end{array}
\]
For $c=[t^{12}]\in C$, the integral maps
\[
 L:C\oplus C\to C,\quad L(b,a)=(1-s)b+s(1-r)a,
 \qquad\widetilde L(b,a,h)=L(b,a)-hc
\]
have Smith forms $(1^{96},2)$ and $(1^{97})$, respectively.
All 24 automorphism certificates have rank $97$ on $C/pC$ for
$p=2,3$. After the exact integral cusp gauge, their two hexagons
hold in $G_6$ and the resulting 24 automorphisms are distinct.
In particular twelve have the extra CM coordinate. No discrepancy
with these numerical assertions of R12 was found.
\end{statement}
\begin{proof}
The independent verifier is \texttt{code/de\_verify\_cm.py}.
It imports none of the source programs. It enumerates determinant-one
matrices modulo twelve and takes their classes by simultaneous sign.
It constructs the graph with edges $Q$ and vertices
$Q/\langle s\rangle\amalg Q/\langle r\rangle$. There are $576$
edges and $480$ vertices. The graph is connected, and the fundamental
cycles for a spanning tree form an integral basis of its homology;
their coordinates on the 97 omitted edges form $I_{97}$.
The Bass--Serre action identifies that homology with $N_B^{\mathrm{ab}}$.
It also identifies $N_B$ with the free group of rank $576-480+1=97$,
because $N_B$ meets both conjugates of the finite factors trivially.

For a word $w\in B$, use its image $q(w)$ and integral path chain
$P(w)$, with
\[
 P(s)=0,\quad P(r)=e_1-e_r,\quad
 P(ab)=P(a)+q(a)P(b).
\]
Chains distinguish the elements of $G_j$ after reduction modulo $j$:
two lifts of the same element of $Q$ differ by a cycle, which is zero
modulo $j$ precisely when their difference lies in $K_j$. Thus all
group equalities tested in this model are equalities in $G_j$.

Let $H=q(F)$. Direct enumeration gives $|H|=96$,
$|C_H(r)|=6$ and $|C_H(s)|=8$.
The verifier lists all 48 pairs $(A,\beta)$ in these two centralizers.
It uses the unique lift of each projective matrix whose reduction
modulo four is in the four-group generated by the reductions of
$u^2,v^2$. For $A^{-1}$ the action on the rational three-torsion line
$\ker(R-I)$ gives one sign. The reduction of $\beta^{-1}$ modulo
three, lying in $\{\pm I,\pm S\}$, gives the other. Equality of the
two signs is the CM-compatible class; their inequality is the extra
class. These are exactly the restrictions to
$M'=\Q(i,(3/4)^{1/4})$ used in LM7 D.25. The program computes
24 pairs of each sign.

Choose integral paths $P_A,P_\beta$ lifting the pair. The cusp error
for $s\mapsto\beta s\beta^{-1}$ and $r\mapsto ArA^{-1}$ is
\[
 E=P_\beta-sP_\beta+
       s(P_A-rP_A+A P(r)-P(r)).
\]
Changing the paths by cycles $b,a$ changes this error by $L(b,a)$.
Exact Gaussian elimination over $\F_2$ and $\F_3$ gives the displayed
solvability counts. The verifier then reads the stored source
certificates as proposed vectors, substitutes them into this equation,
and computes the induced map on all fundamental cycles. The edge at
$1\in Q$ maps to
$P_A-P_\beta+e_A$; translating this formula by $Q$ gives every other
edge. The resulting matrices have the asserted ranks. They act
invertibly on $C/6C$ and identically on $Q$, proving that each proposed
endomorphism of $G_6$ is an automorphism.

To verify the commutator gauge independently, the program projects
chains to the corresponding six-edge graph for $S_3$. Its two
integral basis cycles are the images of $x,y$. Thus it computes the
actual integral abelianization $\beta^{\mathrm{ab}}=ax+by$ of each
chosen word. With $k=b-a$ and $\beta'=x^k\beta$,
\[
 f=s\beta's(\beta')^{-1}\in F',
\]
because conjugation by $s$ exchanges $x,y$ in $F^{\mathrm{ab}}$ and
$\beta'$ has equal coordinates. The program directly multiplies the
words for both $m=0$ hexagons in $G_6$; all 48 equalities hold.
It also compares the resulting images of $s,r$, obtaining 24 distinct
automorphisms. Independently it finds
\[
 q(C_{G_2}(s)\cap\im F)=1,
 \qquad q(C_{G_3}(r)\cap\im F)=1,
\]
by solving the corresponding centralizer equations. These checks
account for the finite conjugator ambiguity in the recorded
coordinates. The gauge replaces $A,\beta,V$ by simultaneous left
translates, so $T=A^{-1}V$ and $W=\beta^{-1}V$ are unchanged.

Finally the verifier computes the integer matrices of $L$ and
$[L,-c]$ in its newly constructed cycle basis. Exact Smith reduction
gives the stated invariant factors; Hermite reduction of $[L,-c]$
gives $I_{97}$. The cycle $c$ has relatively prime coordinates.
The outputs retain every pair, rank, abelianization and gauge, together
with the complete source certificate files. These integer calculations
confirm more than the ranks of the maps modulo two and three.
\end{proof}

\begin{statement}{E.2}{Open}{The full profinite gentle CM image}
We have not determined whether $\widehat{\mathrm{GT}}_0$ reaches the
extra CM coset. E.9 proves that \emph{none} of the twelve particular
extra shadows recorded by LM7 D.25 is genuine, using an explicit
common finite refinement. The full existence question ranges over
every shadow with the specified CM matrices. E.10 computes the
affine parameters in those matrix fibres; the twelve recorded
representatives do not exhaust them. Thus the asserted equivalence
in Task E.2 between the full CM question and genuineness of one of
those twelve representatives still requires a proof of that reduction.
E.11 constructs new points in every extra CM matrix fibre that
survive that same common refinement. E.15 constructs compatible
families in all 24 matrix fibres through the entire relative
class-two tower, with multiplier one. E.16 extends those families
simultaneously through every finite collection of nonabelian simple
quotients, and separately through each finite perfect quotient.
E.20 further extends the class-two families simultaneously through
every relative quotient with normal kernel of order prime to six,
including its arbitrary higher commutator layers. The combined
abelianization and the single cusp gauge are retained. The new
specified comparison outside that family has a relative three-primary
kernel with a nonzero third lower-central term. The unresolved set is the
intersection of the genuine-shadow image with the \emph{whole}
union of extra CM matrix fibres.
\end{statement}

\begin{statement}{E.3}{Open}{The full profinite GT CM image}
We have not determined whether full $\widehat{\mathrm{GT}}$ reaches
the extra CM coset. The finite tests in E.6 impose consequences of
the pentagon established in D.5. E.17 evaluates the complete group
pentagon in the specified pure four-strand quotient of order $5184$.
It excludes fifteen of the twenty-four retained $G_6$ representatives,
but its rank-one affine equation has actual generating solutions in
every whole CM matrix fibre, simultaneously satisfying the class-two
cusp obstructions of E.11. E.18 retains the next central Ferrari
layer in a pure quotient of order $15552$: its pentagon has the
same zero sets on all whole affine fibres, and all twenty-four
retained words still generate the order-$648$ Ferrari image.
E.19 imposes the modulus-nine pentagon together with the complete
three-strand hexagons on the exact induced core. Each retained
point has one compatible correction in the full three-dimensional
commutator correction image, with generation verified in both
targets.
The missing step is compatibility through
all four-strand quotients. No positive or negative conclusion for
Task E.3 follows from the surviving finite candidates.
The primary comparison in
\cite[Proposition~4.4 and Theorem~4.16]{DE_NakamuraTsunogaiYasuda2010}
constructs the triangle-cover Tate action for full
$\widehat{\mathrm{GT}}$ but identifies it with the normalizer
$h$-coordinate on the Galois subgroup. The required comparison
of these two actions on the entire source remains unproved here.
The newer action of
\cite[Theorem~1.1]{DE_LochakNakamuraSchneps2026}
is defined on the profinite mapping class groups of type $(g,0)$
and $(g,1)$; its Section~8 explicitly leaves compatibility with
the indicated cutting homomorphisms unresolved for full
$\widehat{\mathrm{GT}}$. It does not provide the needed CM
comparison in the checked statements.
\end{statement}

\subsection{The homology refinements and an additional finite quotient}

\begin{statement}{E.4}{Proved}{The checked homology lifting mechanism and its scope}
The Smith computation in E.1 gives the successive corrections used in
LM7 D.26. Each of its twelve extra CM pairs admits a compatible family
of shadows through the quotients $G_j$, with a coherently chosen
level-six cusp gauge and multiplier $\lambda\equiv1\pmod{72}$.
The family of kernels $K_j$ fails to be cofinal among the finite-index
normal subgroups of $B$ lying below $K_6$.
\end{statement}
\begin{proof}
The exact path equation for $t\mapsto t^{1+12h}$ is
$E+L(b,a)-hc=0$, because $q(t)c=c$. At $p=2,3$ start from an E.1
solution modulo $p$ and $h=0\bmod p$. Surjectivity of $\widetilde L$
over $\Z$ cancels the divided error successively modulo $p$; this
constructs solutions over $\Z_p$. Their matrices on $C\otimes\Z_p$
remain invertible, since their reductions modulo $p$ are the verified
matrices. They therefore give automorphisms of the extensions of $Q$.

For $p>3$ choose the identity automorphism. The required matrix
conjugators lift to centralizers: a lift of $A\in C_Q(r)$ has error
$z$ with $(1+r+r^2)z=0$, and
$(1-r)(2+r)=3-(1+r+r^2)$ corrects it after dividing by three.
The analogous formula for $s$ divides by two. Both divisions are
valid in $C\otimes\Z_p$. Chinese remainders on $C/jC$ assemble these
automorphisms in
\[
 \mathcal G=\varprojlim_j G_j
   =\widehat B/\overline{[\widehat N_B,\widehat N_B]}.
\]
The equality follows on each fibre above $Q$ by completion of the
integer cycle group $C$. The induced abelianization of the image of
$\widehat F$ is $\widehat\Z^2$, since the displayed kernel lies in
$\widehat F'$. One global cusp gauge thus makes $f$ a commutator
and gives compatible shadows. The congruences on $h$ give the stated
congruence modulo $72$. A discrete representative matching the global
conjugator both in $G_6$ and in $F^{\mathrm{ab}}/36$ gives its level-six
gauge. This specifies the gauge; it does not assert compatibility for
every independent choice of a finite representative.

We reran the original refinement program. Besides the exact integer
normal forms, it checks the powers $2,4,8,16$ and $3,9,27,81$ and
verifies 384 general-multiplier hexagon identities. These finite
checks support the formula; the preceding correction argument supplies
all powers. The quotient in E.5 has kernel contained in $K_6$ and
does not kill $[N_B,N_B]$, whereas every $G_j$ kills that subgroup.
The factorization test for the quotient maps proves the stated
failure of cofinality.
\end{proof}

\begin{statement}{E.5}{Proved}{All certified shadows survive an explicit $A_5$ refinement}
Define $a:B\twoheadrightarrow A_5$ by
$a(s)=(12)(34)$ and $a(r)=(135)$, and put
\[
 J_{A_5}=K_6\cap\ker a.
\]
The diagonal map induces an isomorphism
$B/J_{A_5}\simeq G_6\times A_5$. Its kernel does not contain
$[N_B,N_B]$. Every one of E.1's 24 certified shadows extends to a
shadow at the inverse image of $J_{A_5}$ in $B_3$. All twelve extra
CM coordinates survive this finite test.
\end{statement}
\begin{proof}
The independent program \texttt{de\_a5\_refinement.py} enumerates
the permutation group generated by the specified two elements.
It has order 60, is perfect and simple, and its cusp $a(t)$ has
order five. The program also verifies the derived-series orders
$576,96,8,1$ for $Q$. Thus $G_6$, an extension of $Q$ by the
abelian group $C/6C$, is solvable. The image of the diagonal map
is a subdirect product. Goursat's lemma expresses a proper subdirect
product by a common nontrivial quotient of the two factors. Every
quotient of $G_6$ is solvable, whereas the only nontrivial quotient
of $A_5$ is $A_5$ itself and is perfect. Hence the image is the
whole product.

The subgroup $a(N_B)$ is normal in $A_5$, and
$A_5/a(N_B)$ is a quotient of $Q$. The same argument gives
$a(N_B)=A_5$, and therefore $a([N_B,N_B])=A_5$.
This proves $[N_B,N_B]\not\subset J_{A_5}$.

Let $\alpha_6$ and $f_6$ be any certified, gauged automorphism and
commutator from E.1. On the product choose
$(\alpha_6,\mathrm{id}_{A_5})$, multiplier $\lambda=1$, and
$f=(f_6,1)$. The image of $F$ is
$\im(F\to G_6)\times A_5$: its projection to $A_5$ is surjective
because its quotient has order dividing six, and the same Goursat
argument applies to its solvable first factor. Taking derived groups
shows that $f$ is the image of an element of $F'$. Both hexagons hold
componentwise, and the product automorphism verifies generation.
This proves the shadow claim by the finite definition.

The refined cusp has order $\mathrm{lcm}(72,5)=360$ and the shadow
order parameter is $180$; $m=0$ and $\lambda=1$ satisfy its unit
condition. The projection $G_6\times A_5\to G_6$ gives the exact
reduction map that retains the original CM coordinate. This quotient
is additional to, and has not been identified with, LM7 D.26's
$S_3$-core quotient $B/J$.
\end{proof}

\subsection{Finite tests from the full GT distribution theorem}

\begin{statement}{E.6}{Proved}{Five simultaneous finite cyclic-cover tests preserve all 24 candidates}
Let $f_0\in F'$ be the explicit gauged word of any E.1 certificate,
and let $q_F:F\twoheadrightarrow H=q(F)$ be the matrix quotient.
There is one word $k\in[N_B,N_B]$ such that $f_*=f_0k$ has the
same image as $f_0$ in $G_6$ and, simultaneously for all
$N\in\{2,3,4,6,12\}$, satisfies
\[
 q_F\delta_N(f_*)=q_F(y)^{c_N(f_*)}q_F(f_*),
 \qquad c_N(f_*)\text{ read modulo }6.
\]
It can be chosen also to satisfy, modulo six,
$c_4=2c_2$, $c_6=c_2+c_3$, and $c_{12}=2c_2+c_3$.
These are assertions in the listed finite targets.
\end{statement}
\begin{proof}
The programs \texttt{de\_distribution\_test.py} and
\texttt{de\_distribution\_more.py} reconstruct words for the 97
fundamental cycles of the graph used in E.1, rewrite them in $x,y$,
and evaluate the maps below. Put $A=H\times\Z/6$, and define
\[
 \psi_N:H_N\cap F\to A,\qquad
 w\longmapsto(q_F\delta_N(w),c_N(w)\bmod6).
\]
The $x$ exponent of every member of $N_B$ is even. Consequently
$K_6\subset H_N$ for every $N\mid12$; also $N_B'\subset F'$.
Induce $\psi_N$ to a homomorphism from $F$ into
$A^N\rtimes\Z/N$ by the transversal $1,x,\ldots,x^{N-1}$.
Its coordinates are the Schreier words
$x^iwx^{-(i+s)}$, with residues taken modulo $N$ and
$s$ the $x$ exponent of $w$ modulo $N$.
If an element is written $(v_i,s)$, multiplication is
$v_i(w_1w_2)=v_i(w_1)v_{i+s(w_1)}(w_2)$.

The normal subgroup generated by commutators of the 97 images is
the image of $N_B'$. Every such commutator has shift zero.
Conjugation of a shift-zero vector by $(v_i,s)$ sends its $i$th
coordinate to $v_i k_{i+s}v_i^{-1}$. Starting from the commutators,
the program repeatedly takes coordinate subgroup closures and these
conjugations until stable. This computes the coordinate projections
exactly: every generator added is a coordinate of an actual product
of conjugates, and the final projections contain all such products.
The coordinate-zero orders obtained are
\[
\begin{array}{c|rrrrr}
 N&2&3&4&6&12\\\hline
 |\psi_N(N_B')|&8&48&48&48&48\\
 |\psi_N(K_6)|&8&144&96&144&288\\
 \mathrm{passing\ compatible\ candidates}&12&12&12&12&12\\
 \mathrm{passing\ extra\ candidates}&12&12&12&12&12.
\end{array}
\]
For each certificate it intersects the exact coset
$\psi_N(f_0)\psi_N(N_B')$ with
\[
 \{(q_F(y)^c q_F(f_0),c):c\in\Z/6\}.
\]
Every intersection is nonempty. The order of $q_F(y)$ is six,
so the exponent modulo six suffices for this target.

To impose these tests simultaneously, \texttt{de\_distribution\_joint.py}
uses a faithful degree-17 permutation representation of $A$:
$H\simeq C_2^2\times\mathrm{SL}_2(\F_3)$ acts regularly on four
points and naturally on the eight nonzero vectors of $\F_3^2$,
while $\Z/6$ acts on disjoint sets of two and three points.
The product of the five induced representations has degree
$17(2+3+4+6+12)=459$. The program takes the derived subgroup of
the image of the 97 $N_B$ generators and projects its shift-zero
elements to coordinate zero in each factor. Exact permutation-group
membership gives the joint image of $N_B'$ in $A^5$, of order
$42\,467\,328$.

For each certificate we test the finite target tuples with
$(c_2,c_3,c_4,c_6,c_{12})=(a,b,2a,a+b,2a+b)$ modulo six.
All 24 required coset intersections are nonempty. A tuple in this
joint image is attained by a single word of $N_B'$, which gives
the asserted common correction. The complete output includes the
chosen coefficient tuple, the derived-image generators, its order,
and the membership results. Thus no inference from five independent
choices is used for the simultaneous statement.
\end{proof}

These computations concern the explicit projections $\psi_N$.
They do not enumerate all shadows at a refinement carrying these
projections, and they do not impose generation there. In particular
they show that these necessary distribution projections do not
exclude an extra CM candidate. D.6 still requires the full collection
of finite shadow fibres. The theorem in D.5 provides a concrete
additional relation to test, but E.2 and E.3 remain the precise
unsettled questions.

\subsection{The requested \texorpdfstring{$S_3$}{S3}-core refinement}

\begin{statement}{E.7}{Proved}{All 24 candidates lift through the specified $S_3$ core}
Fix the free basis $n_1,\ldots,n_{97}$ of $N_B$ given by the
fundamental cycles in E.1, ordered by the omitted-edge order in its
program. Let $\rho:N_B\twoheadrightarrow S_3$ send $n_1,n_2$ to
two different transpositions and the other basis elements to $1$.
Put
\[
 J=\operatorname{core}_B(\ker\rho\cap K_6),
 \qquad D_J=B/J.
\]
There is an exact sequence of finite groups
\[
 1\longrightarrow (\Z/3)^{288}\longrightarrow D_J
       \longrightarrow G_6\longrightarrow1.
\]
The cusp $t$ has order $72$ in $D_J$. Each of E.1's 24 certified
shadows is the reduction of a shadow at the inverse image of $J$ in
$B_3$, with multiplier one. All twelve extra CM shadows consequently
survive this first requested nonabelian refinement.
\end{statement}
\begin{proof}
The complete finite calculation is \texttt{de\_s3\_core.py}; it was
run to completion. We describe its faithful model and the exact
surjectivity test. Identify $S_3$ with the affine group of $\F_3$,
whose elements are pairs $(\epsilon,b)$ acting by
$z\mapsto\epsilon z+b$, with $\epsilon\in\{1,-1\}$.
The two transpositions are $(-1,0)$ and $(-1,1)$.
On the graph of E.1 label the first and second omitted edges by
these elements, and every other edge by $1$. Denote these labels
by $a_q$ for $q\in Q$. The induced homomorphism
\[
 \omega:B\longrightarrow S_3^Q\rtimes Q
\]
has $s$ component identically $1$, right permutation $q\mapsto qs$,
and $r$ component $a_q a_{qr}^{-1}$ with right permutation
$q\mapsto qr$. The identities $s^2=r^3=1$ follow by cancellation.
For $n\in N_B$, the component at $1$ is $\rho(n)$; the other
components are its conjugate Schreier representations. Hence
$\ker\omega$ on $N_B$ is the normal core of $\ker\rho$.
The combined map $(B\to G_6,\omega)$ has kernel exactly $J$ and
is a faithful model of $D_J$.

Since every element of $S_3$ has sixth power one,
$\omega(K_6)=\omega(N_B')$. Each of its components lies in
$A_3\simeq\F_3$. This image is therefore a vector subspace
$M\subset\F_3^Q$. The program forms the commutators of the images
of the 97 basis elements and closes their span under conjugation
by $s,r$. Conjugation by $(\epsilon_q,b_q;g)$ is the linear map
\[
 (z_q)_{q\in Q}\longmapsto
       (\epsilon_q z_{qg})_{q\in Q}.
\]
The iteration stabilizes at dimension $288$. Every inserted vector
is the image of a product of conjugates of the basis commutators;
conversely a stable span contains the image of their normal closure.
Thus the computed subspace is exactly $\omega(N_B')$.
The kernel of $D_J\to G_6$ is $M$ by the combined faithful model.
Direct evaluation gives $t^{72}=1$; its order is at least 72 from
its image in $G_6$.

For each original pair choose the already recorded words $A,\beta$
and write $S'=\beta s\beta^{-1}$, $R'=ArA^{-1}$ in this model.
The error $S'R't^{-1}$ lies in $M$. Replacing the conjugators by
$b\beta,aA$, with $b,a\in M$, changes the error by
\[
 (1-S')b+S'(1-R')a.
\]
The conjugation actions in this equation are the just specified
linear maps on $M$. For every one of the 24 pairs the matrix of
this map has rank 288. The program solves the equation and records
its 576 correction coordinates. Each correction is represented by
an actual word of $N_B'$, by the construction of $M$.

The proposed images of $s,r$ satisfy their finite-order relations,
so define a homomorphism $B\to D_J$, whose composition with
$D_J\to G_6$ is the original automorphism. To check surjectivity
on the kernel, evaluate all 97 $n_i$ on the proposed images, take
their pairwise commutators, and close their vector span under the
two proposed generator actions. Their sixth powers are all trivial
in the wreath component. These operations give exactly the image
of $K_6$. In every case the resulting span has dimension 288.
The map is therefore onto $D_J$: it is onto $G_6$ and onto the
entire kernel. The first deterministic solution already has this
property for all 24 pairs; no probabilistic inference is used.

Finally the correction words lie in $N_B'\subset F'$, so the
integral abelianization of $\beta$ and the cusp gauge from E.1 are
unchanged. After that gauge the program directly verifies both
hexagons in the wreath component; they hold in $G_6$ by E.1 and
hence in the faithful combined model. The commutator condition
holds for the same abelianization reason. With $m=0$ and
$\lambda=1$, the unit condition and the verified surjectivity give
a finite GT shadow. Reduction by $D_J\to G_6$ is the specified
original shadow. This proves the claim for all 24 recorded fibres.
\end{proof}

For reproducibility, the first basis word in this instance is
$(sr^{-1})^{12}$; the complete first two words and the deterministic
basis construction are in the JSON and source. LM7 D.26 left the
free basis in this construction to be chosen. We fix it here rather
than identifying cores arising from different choices. The next
calculation excludes all twelve representatives that survived this
particular core.

\subsection{A finite exclusion of all twelve recorded extra shadows}

\begin{statement}{E.8}{Known}{Finite-factor conjugacy for the full profinite gentle action}
Every $g=(\lambda,f)\in\widehat{\mathrm{GT}}_0$ induces a continuous
automorphism of $\widehat B=C_2\amalg C_3$ sending $s$ to a conjugate
of $s$ and $r$ to a conjugate of $r^{\lambda\bmod3}$.
For $\lambda\equiv1\pmod3$, the conjugators can be chosen in
$\widehat F=\ker(\widehat B\to S_3)$.
\end{statement}
\begin{proof}
Harbater--Schneps, Proposition 4, proof of $(i)\Rightarrow(iii)$,
explicitly extends the full profinite gentle action by
$u\mapsto u^\lambda$, $v\mapsto f^{-1}v^\lambda f$ on
$\widehat B_3$, and proves that its central generator goes to its
$\lambda$th power \cite{DE_HarbaterSchneps2000}.
It descends to the quotient by the closed center. The universal
property of profinite completion identifies this quotient with the
free profinite product $C_2\amalg C_3$: continuous homomorphisms
from either to a finite group are precisely the pairs of elements
satisfying $s^2=r^3=1$.

Scheiderer's Appendix to Lochak--Schneps, unnumbered Corollary on
printed p.598, places every nontrivial finite subgroup of a free
profinite product in a conjugate of a factor
\cite{DE_LochakSchneps1997}. Applying it to the images of
$\langle s\rangle$ and $\langle r\rangle$ gives the stated conjugacy
possibilities. The induced automorphism of
$B^{\mathrm{ab}}=C_2\times C_3$ is multiplication by $\lambda$,
as follows from the abelianization of the displayed braid action.
This distinguishes $r$ from $r^{-1}$ and selects the stated power.
For $\lambda\equiv1\pmod3$, reduction of the braid action to $S_3$
fixes both braid generators, because $\lambda$ is odd and
$f\in\widehat F'$. A conjugator for $s$ consequently projects into
$C_{S_3}(s)=\langle s\rangle$, and a conjugator for $r$ projects
into $C_{S_3}(r)=\langle r\rangle$. Right multiplication by the
appropriate power of $s$ or $r$ leaves the conjugate unchanged and
makes that projection trivial. This gives conjugators in $\widehat F$.
\end{proof}

Put $N=N_B$ for the rest of this subsection, and define characteristic
subgroups of $N$, hence normal subgroups of $B$, by
\[
 J_3=(\gamma_3N\,N^3)\cap(N'N^2),\qquad
 J_2=(\gamma_3N\,N^4)\cap(N'N^3),\qquad
 J_*=J_2\cap J_3.
\]
All are contained in $K_6=N'N^6$. Write $D_i=B/J_i$ and let
$\widetilde J_i$ be the inverse image of $J_i$ under $B_3\twoheadrightarrow B$.
These are finite-index normal subgroups of $B_3$ contained in $PB_3$.
The projections $D_*\to D_i\to G_6$ induce the reduction maps on
gentle shadows used below. Their pure parts surject onto
$H=q(F)\simeq C_2^2\times\mathrm{SL}_2(\F_3)$ by restriction of $q$.

\begin{statement}{E.9}{Proved}{An explicit common refinement removes all twelve recorded extra shadows}
Number the 24 E.1 certificates from zero in their stored source order,
retaining exactly the records containing \texttt{mod6\_beta\_cycles}.
The exact gentle-shadow fibres of the two projections have the following
nonempty instances among these certificates:
\[
\begin{array}{c|l|rr}
\text{refinement}&\text{surviving certificate indices}
 &\text{compatible}&\text{extra}\\\hline
 D_3\to G_6&0,1,2,3&2&2\\
 D_2\to G_6&0,16&1&1\\
 D_*\to G_6&0&1&0.
\end{array}
\]
The extra indices are
$1,2,4,7,9,10,12,15,16,19,21,22$.
Consequently none of the twelve recorded extra shadows is in the
image of $\widehat{\mathrm{GT}}_0$, and none is in the image of
$\widehat{\mathrm{GT}}$. This excludes those specified finite
shadows, while E.2 asks about all shadows in their CM matrix fibres.
\end{statement}
\begin{proof}
We give the exact linear equations, including the multiplier and
commutator conditions. Since $N$ is free of rank $n=97$, its image
in $D_3$ is the product of the free class-two exponent-three group
on $n$ generators with $C/2C$. The kernel of $D_3\to G_6$ is
\[
 M_3=\Lambda^2(\F_3^{97}),\qquad \dim M_3=4656.
\]
For $D_2$ the corresponding product is the free class-two
exponent-four group on $n$ generators with $C/3C$. Squares and
commutators have order two and generate the central kernel
\[
 M_2\simeq(\Z/2)^{97+\binom{97}{2}},\qquad\dim M_2=4753.
\]
These descriptions follow directly from the normal forms
\[
 (a,z)(b,w)=(a+b,z+w+\tfrac12a\wedge b)\quad(p=3)
\]
and
$n_1^{a_1}\cdots n_{97}^{a_{97}}\prod_{i<j}[n_i,n_j]^{z_{ij}}$
with $a_i\bmod4$, $z_{ij}\bmod2$ at two. Collecting these normal
forms verifies the defining exponent and class-two relations, and
their independent generator and commutator coordinates prove that
no further relations have been imposed. The two prime parts form
the stated products since their orders are coprime.

Let $\mathcal S,\mathcal R$ denote conjugation by $s,r$ on $M_i$.
At three these are the exterior-square actions induced from $C/3C$.
At two their duals act on the quadratic functions on $C/2C$ by
$q(v)\mapsto q(Sv)$ and $q(v)\mapsto q(Rv)$; here $S,R$ are the
verified 97-dimensional cycle actions. This formula follows by
squaring the collected normal form, whose square is the quadratic
function of its parity vector. It includes both square and
commutator coordinates.

For one recorded certificate, undo its explicit E.1 cusp gauge and
take its stored word $\beta$. Set
\[
 S_0=\beta s\beta^{-1},\qquad R_0(\lambda)=S_0t^\lambda.
\]
The first element is an involution. For every possible lift of the
same finite shadow, the new images have the unique form
\[
 S'=S_0k,\qquad R'=k^{-1}R_0(\lambda),\qquad k\in M_i.
\]
The action of $S_0$ on $M_i$ is $\mathcal S$, and that of
$R_0(\lambda)$ is $\mathcal R$: their images in $Q$ are $s,r$,
and $N$ acts trivially on $M_i$.
The exact order equations are therefore
\begin{equation}\label{eq:DEgentlenorm}
 (1+\mathcal S)k=0,\qquad
 (1+\mathcal R+\mathcal R^2)k=R_0(\lambda)^3.
\end{equation}
We explain why these test all finite gentle lifts, without assuming
that torsion in a finite group is conjugate to a specified factor.
Put $f=sS'$ after undoing the gauge, $u=t^{-1}$, and $v=sus$.
For the two finite hexagons, multiplying the first on the right by
$u^\lambda$ gives $f^{-1}s$, while multiplying the second on the
left by $u^\lambda$ gives $sf$. Their left sides are identical.
Thus they imply $(sf)^2=1$. After this equality the first
hexagon reads $(S't^\lambda)^3=1$. Reversing these algebraic steps
proves the converse. In particular \eqref{eq:DEgentlenorm} is the
complete pair of hexagon equations at the indicated quotient.
The gauge is conjugation by a fixed power of $x=t^{-2}$ and
commutes with $t^\lambda$, so undoing it and reinstating it are
inverse operations throughout this calculation.

In $D_3$ the cusp has order 72. Its order is at least 72 from $G_6$
and divides 72 from the exponent-three and exponent-two factors.
The finite order parameter is therefore 36, and the multiplier
over a recorded shadow has the unique residue $\lambda=1\bmod72$.
Every $k\in M_3$ is in the image of $F'$. As two is invertible,
the first equation in \eqref{eq:DEgentlenorm} is parametrized by
$k=(1-\mathcal S)w$. The program
\texttt{de\_class2\_gentle.py} constructs the complete matrix
\[
 (1+\mathcal R+\mathcal R^2)(1-\mathcal S):M_3\to M_3
\]
and all 24 right sides by word collection. Its rank is 1542.
Exact elimination of the augmented matrix, whose rank is 1547,
gives exactly indices $0,1,2,3$. Each retained correction is
substituted into both original equations.

In $D_2$ the cusp has order 144. Indeed $t^{12}$ has the primitive
cycle from E.1, so its sixth power has a nonzero square coordinate
modulo two, while its twelfth power is trivial in both factors.
The finite order parameter is 72. All multiplier lifts are
$\lambda=1,73\bmod144$; both are tested. For the commutator
condition let $L_N$ be the image of $N$ in $F^{\mathrm{ab}}=\Z^2$.
Exact word abelianization gives the column lattice
\[
 L_N=\begin{pmatrix}6&4\\0&2\end{pmatrix}\Z^2.
\]
The image of $J_2$ in this abelianization is exactly $12L_N$:
membership gives divisibility by four and three, and twelfth powers
of elements of $N$ give the reverse inclusion. A square coordinate
in $M_2$ is represented by $n_i^6$ and a commutator coordinate by
a commutator. Hence $k\in M_2$ lies in the image of $F'$ exactly
when the two coordinates of its square contribution to $6L_N/12L_N$
vanish. The program \texttt{de\_class2\_two\_gentle.py} appends
these two rows to \eqref{eq:DEgentlenorm}. The coefficient matrix
has rank 3927; its augmentation by all 48 right sides has rank 3932.
Its complete solution test gives exactly indices $0,16$.
Every retained solution is substituted into the two order equations
and the two abelianization rows. The fixed E.1 gauge made the
original $f$ a commutator, so this is precisely the condition that
the corrected, regauged $f$ be a commutator as well.

All passing solutions satisfy generation. The homomorphism
$B\to D_i$ defined by $S',R'$ preserves the image of $N$, because
its map to $Q$ is the recorded identity before gauge. The defining
class and power laws show that it kills $J_i$. On each prime part
of the image of $N$, its Frattini quotient map is exactly the
invertible E.1 map on $C/pC$. The elementary Frattini argument for
a finite $p$-group therefore makes it surjective on both prime
parts, and hence on $D_i$. The multiplier, commutator, hexagon and
generation conditions of a finite gentle shadow have all been
checked.

The bit-packed elimination routines in \texttt{de\_gf3.c} are
independently compared with the earlier Python elimination on four
rectangular matrices for each prime, including multiple machine
word boundaries. An additional independent calculation
\texttt{de\_class2\_direct\_check.py} uses the full 9312 by 4656
matrix of the peripheral correction map at three, agrees with the
separate Jordan-block computation, and gives the same four
survivors after all conjugator-centralizer freedoms. The full
hexagon calculation above supplies the exact finite-shadow scope.

Every shadow over $\widetilde J_*$ reduces to shadows over both
$\widetilde J_2$ and $\widetilde J_3$. Intersecting the two
displayed index lists leaves only index zero. That record has
$A=\beta=1$ and zero cycle corrections; its identity shadow lifts
to every quotient, including $D_*$. Thus the common-refinement
image among the 24 recorded points is exactly the claimed singleton.
The all-finite-quotient criterion D.6 excludes the other 23 from
the full profinite gentle image. The inclusion
$\widehat{\mathrm{GT}}\subset\widehat{\mathrm{GT}}_0$ gives the
stated full GT exclusion as well.
\end{proof}

\begin{statement}{E.10}{Proved}{The affine sizes of the CM matrix fibres}
For each of the 24 admissible CM matrix pairs and multiplier one
modulo 72, the cusp-preserving changes of its two homology
conjugators form affine spaces of dimensions 98 over $\F_2$ and
97 over $\F_3$. Quotienting by the changes that leave both generator
images fixed gives dimensions 17 and 18. Quotienting further by
simultaneous conjugation by cusp powers gives dimensions 16 and 17.
The finite classification modulo all inner conjugations by the
centralizer of $t$ in $C/pC$ has dimensions 8 and 9.
Invertibility of the resulting action on $C/pC$ is an additional
condition in each of these spaces.
\end{statement}
\begin{proof}
The exact program \texttt{de\_affine\_fibres.py} computes the
kernel of $L$ from E.1 and the following subspaces. Over $\F_2$
the dimensions of $C^s,C^r,C^t,C^Q$ are respectively $52,29,9,0$;
over $\F_3$ they are $45,34,9,0$. The changes that fix both
generator images are $C^s\oplus C^r\subset\ker L$. Simultaneous
inner conjugation adds the diagonal copy of $C^t$; its intersection
with that first subspace is the diagonal $C^Q=0$. The cusp powers
give the one-dimensional diagonal subspace generated by the
primitive cycle $c=[t^{12}]$. Taking the indicated dimensions gives
all assertions. The output stores explicit bases and complements,
and directly substitutes every column into $L$.

Conjugation by an actual power of $t$ centralizes $t$ in every
quotient, so that operation can be transported through every
refinement. The larger finite space $C^t$ is classified here only
at the homology quotient: transporting one of its elements as a
centralizer would require a lift that centralizes $t$. Consequently
the last quotient does not by itself replace the whole finer
lifting problem. In particular the twelve original representatives
cannot be treated as exhaustive matrix fibres. R11 supplies a
genuine arithmetic point in each compatible CM matrix fibre, while
E.9 excludes the recorded points in every compatible fibre except
the identity fibre. These give concrete distinct points with the
same CM matrices.
\end{proof}

\begin{statement}{E.11}{Proved}{Every admissible CM matrix fibre has a replacement at the common refinement}
Each of the 24 admissible CM matrix pairs of E.1 is attained by a
gentle shadow at $\widetilde J_*$, where $J_*=J_2\cap J_3$ is the
explicit refinement of E.9. In particular all twelve extra CM
\emph{matrix fibres} have nonempty image from this refinement,
although their twelve original recorded representatives are excluded
by E.9. The new witnesses are finite shadows; their membership in
the full profinite image is still to be determined.
\end{statement}
\begin{proof}
The programs \texttt{de\_fibre\_three.py} and
\texttt{de\_fibre\_two.py} use the complete 18- and 17-dimensional
spaces after conjugator centralizers in E.10. They construct exact
five-component quadratic obstruction maps on these spaces. The
three-primary map projects the eight alternating coinvariants
modulo their three centralizer directions. The two-primary map
projects the sixteen quadratic coinvariants modulo their eleven
centralizer and multiplier directions. Both maps are checked
against direct collection of words for every matrix pair. Searching
their finite parameter spaces gives a zero in each of the 24
fibres, and the induced action has rank 97 in every case. Only
these produced and verified witnesses are used; the search is not
used to infer any negative assertion.

Write the two pairs of changes of the conjugator cycles as
$(b_2,a_2)$ and $(b_3,a_3)$. The program
\texttt{de\_fibre\_crt\_verify.py} uses
\[
 b=3b_2+4b_3\pmod6,\qquad
 a=3a_2+4a_3\pmod6,
\]
and prefixes the original words $\beta,A$ by
$\prod n_i^{b_i},\prod n_i^{a_i}$. These are literal words in
$s,r$ stored in its output. Their graph cusp equation and both
97-dimensional generation ranks are checked again. Their images
in $Q$ are the original matrices, so their CM coordinates agree
with the prescribed pair, also after the common cusp gauge.

The complete norm systems of E.9 are then solved anew for these
CRT words. Every one of the 24 points has a solution at both
primes with $\lambda=1\bmod144$. The two-primary system includes
the two exact commutator-image rows. Let $k_2$ be its square and
commutator coordinates and $k_3$ the three-primary commutator
coordinates. The common correction is the explicit word
\[
 k=\prod_i n_i^{6(k_2)_i}
       \prod_{i<j}[n_i,n_j]^{3(k_2)_{ij}+4(k_3)_{ij}}.
\]
Its two-primary component is $k_2$ and its three-primary component
is $k_3$. For the integral cusp gauge $g$ computed from the new
word $\beta$, set
\[
 f=sx^g(\beta s\beta^{-1})k x^{-g}j.
\]
Here $j$ is an explicitly recorded product of twelfth powers of
two basis words of $N$, chosen to cancel the abelianization of $k$.
The two commutator rows ensure that this abelianization belongs to
$12L_N$; an exact unimodular pair in the displayed basis of $L_N$
constructs $j$. Thus $j\in J_*$ and the displayed word $f$ lies
in $F'$ itself. This gives one common commutator word, with a
single integral gauge, for the two primary components.

Finally both hexagons are substituted directly, as words, into a
faithful collector for the product of the exponent-four and
exponent-three class-two images of $N$. There are 48 successful
hexagon checks. For efficiency a factored central block is
evaluated through its exact action under the current $Q$ coset:
at three its alternating matrix changes by $Z\mapsto AZA^t$;
at two a square vector $u$ and symmetric commutator matrix $Z$
change to $Au$ and the upper entries of
$A(Z+\operatorname{diag}u)A^t$. These formulas follow by collecting
the explicit square and commutator factors above, so this evaluation
is precisely substitution of those actual words. The other letters
are collected individually by the omitted-edge basis. A word is
trivial in $D_*$ exactly when its $Q$ image and both collected
primary normal forms vanish.

The checked Frattini ranks give generation on both primary kernels,
as in E.9. The word, multiplier, commutator and generation conditions
therefore give a shadow at $\widetilde J_*$ in every asserted fibre.
The JSON retains the literal CRT conjugators, all factored correction
words, both ranks, the common gauge and multiplier, and every
hexagon result.
\end{proof}

\begin{statement}{E.12}{Proved}{All admissible CM pairs lift through the whole class-two three-primary tower}
For every admissible CM matrix pair in E.1 there is a compatible
family of gentle shadows through all quotients
\[
 B/\bigl((\gamma_3N\,N^{3^k})\cap(N'N^2)\bigr),\qquad k\geq1,
\]
with multiplier one and that CM pair. Thus all twelve extra matrix
fibres survive this entire tower. This assertion concerns matrix
fibres with a coherent cusp gauge; it does not assert that an
independently chosen normalized finite representative extends.
\end{statement}
\begin{proof}
Let $\mathcal P$ be the free class-two pro-three group on 97
generators. Its logarithmic coordinates form
$V\oplus\Lambda^2V$ for $V=\Z_3^{97}$, with product
$(a,z)(b,w)=(a+b,z+w+\frac12a\wedge b)$. The relevant completed
kernel of the extension of $Q$ is
$\mathcal P\times(C/2C)$. The odd denominator two is a unit in
$\Z_3$. Consequently the cusp equation in two conjugators has
polynomial coordinate functions over $\Z_3$, of degree at most two.

We verify that its full derivative is onto at each three-primary
witness of E.11. Its homology part is $L$, of rank 97 modulo three.
The derivative in the central coordinates has rank 4648, with the
eight-dimensional alternating coinvariant quotient computed in
E.9. Changes by the order-three conjugator centralizer give rank
three in this quotient. The remaining five components are exactly
the quadratic map used in \texttt{de\_fibre\_three.py}.
Here is a coordinate description of that map and its derivative.
For an invariant alternating form $B$, write $\tau_{S_0},\tau_t$
for the central parts of conjugation by $S_0=\beta s\beta^{-1}$
and by $t$, evaluated by $B$. A homology change $(b,a)\in\ker L$
changes the central cusp coordinate by
\[
 \tau_{S_0}(a-b)-\tau_t(a)
       +\tfrac12 B\bigl(b,(T-1)a\bigr),\qquad T=SR.
\]
This follows by multiplying
$n(b)\,S_0(n(b)^{-1}n(a))S_0^{-1}$, the old cusp error, and
$t n(a)^{-1}t^{-1}$ in the displayed group law. The central
coordinates chosen for $n(a),n(b)$ cancel in the coinvariant
quotient. The program checks the formula against direct word
collection.

The independent derivative calculation
\texttt{de\_fibre\_three\_smooth.py} evaluates this full
eight-component derivative at each of the 24 zeros, adjoins the
three centralizer directions, and obtains rank eight in every
case. Thus the full coordinate derivative has rank
$97+4648+8=4753$, equal to the dimension of the cusp target.
The centralizer corrections and central corrections used to make
the initial error zero do not change that rank: central corrections
do not change the homology coordinates, and the five-component
quadratic map is constant along the discarded centralizer
directions. All these equalities are checked over $\F_3$.

Successive lifting is now explicit. At a solution modulo $3^k$,
divide its error by $3^k$, reduce modulo three, and solve the
surjective derivative equation for a correction of the two
conjugators by $3^k$ times their coordinates. The polynomial
expansion shows that the corrected error is zero modulo
$3^{k+1}$; its terms of degree at least two in the correction are
divisible by $3^{2k}$. The corrections preserve all preceding
levels, and the derivative retains the same reduction modulo
three. Induction gives convergent conjugators over $\Z_3$.
Their component in $C/2C$ is kept fixed throughout.

The resulting map preserves $Q$ and is invertible on the
three-primary Frattini quotient and on $C/2C$, by the verified
ranks in E.11. The pro-three Frattini argument makes it onto
$\mathcal P$, and its induced finite maps are automorphisms.
Choose profinite lifts of the limiting conjugators to $\widehat F$
and one cusp gauge using their actual abelianizations. The formula
$f=s\beta's(\beta')^{-1}$ then gives a commutator and compatible
hexagons at every finite level. The common multiplier is one.
The simultaneous gauge leaves the CM coordinates unchanged.

This tower kills $\gamma_3N$. The map to $A_5$ in E.5 sends $N$
onto $A_5$, and hence also sends $\gamma_3N$ onto $A_5$.
Its refinement therefore does not factor through any member of
the displayed tower. This explicit factorization obstruction
records why these compatible class-two lifts do not complete E.2.
\end{proof}

\begin{statement}{E.13}{Proved}{Every CM matrix fibre survives the common exponent-eight refinement}
Put
\[
 J_{8,3}=(\gamma_3N\,N^8)\cap(\gamma_3N\,N^3).
\]
This is a normal finite-index subgroup of $B$, contained in $J_*$
of E.9. Its inverse image in $B_3$ is normal, has finite index,
and is contained in $PB_3$. Each of the 24 admissible CM matrix
pairs of E.1 is attained by a gentle shadow at this refinement,
with multiplier one modulo $288$. In particular every one of
the twelve extra CM matrix fibres remains nonempty here.
\end{statement}
\begin{proof}
We compute entire central correction fibres above the 24 affine
homology points of E.11. The image of $N$ in the two-primary
factor is the free class-two exponent-eight group, with ordered
normal form
\[
 \prod_i n_i^{a_i}\prod_{i<j}[n_i,n_j]^{z_{ij}},
 \qquad a_i\in\mathbb Z/8,\quad z_{ij}\in\mathbb Z/4.
 \tag{E.13.1}
\]
The product in these coordinates is
\[
 (a,z)(b,w)=(a+b,z+w-(a_jb_i)_{i<j}).
 \tag{E.13.2}
\]
Indeed moving $n_i^{b_i}$ past $n_j^{a_j}$ for $j>i$
introduces $[n_i,n_j]^{-a_jb_i}$. Commutators have eighth
power one, and the eighth power of a product gives their
twenty-eighth power one. Since $\gcd(8,28)=4$, their fourth
powers vanish. Conversely the ordered normal form with
these relations has exponent eight. This proves the stated free quotient
and the faithful collector used below.

The kernel above the fixed $G_6$ homology image, retaining
$N^{\rm ab}/3$, has a particularly useful description:
\[
 A\cong(\mathbb Z/4)^{4753},\qquad
 e_i=n_i^6,
 \quad e_{ij}=[n_i,n_j]^3.
 \tag{E.13.3}
\]
The two-primary abelian coordinates of an element of this
kernel are even, and multiplication by three is invertible
modulo four, so these elements span the whole kernel.
They are independent by (E.13.1). Moreover their commutators
are $[n_i,n_j]^{36}=1$, so the displayed group is abelian.
This is the exact reason that this entire finite correction
problem is linear over $\mathbb Z/4$.

Let $\beta$ be one of the literal affine conjugators from
E.11, before its cusp gauge, and put $S_0=\beta s\beta^{-1}$,
$R_0=S_0t$. Every lift of the fixed homology shadow has
\[
 S=S_0k,\qquad
 R=k^{-1}S_0t^{1+72h},
 \quad k\in A,\quad h\in\mathbb Z/4.
\]
The complete equations, with the two commutator-image rows,
are consequently one linear system over $\mathbb Z/4$:
\[
 (1+S_0)k=0,\qquad
 (1+R_0+R_0^2)k=(S_0t^{1+72h})^3,
 \qquad B_N^{-1}\operatorname{ab}_F(k)=0\pmod {24}.
 \tag{E.13.4}
\]
Here $B_N=\left(\begin{smallmatrix}6&4\\0&2\end{smallmatrix}\right)$
is the chosen basis matrix for the lattice $L_N$ in E.9.
Define the integer columns $H_i=B_N^{-1}\operatorname{ab}_F(n_i)$.
Then
the two equations are $H(k_i)_i=0\pmod4$, because
$\operatorname{ab}_F(k)=6B_NH(k_i)_i$ and the new kernel
has abelian image $24L_N$. The right side of the second
equation is affine in $h$: $t^{72}\in A$, and $A$ is abelian.
The common multiplier range is $1+72h\pmod{288}$.

For clarity, the action matrices in (E.13.4) use the
ungauged $S_0,R_0$. The normalized pair is obtained by
conjugation by $x^g=t^{-2g}$, which fixes $t$. Thus it
preserves both order equations, and its conjugation of
$k\in F$ preserves the exact $F^{\rm ab}$ coordinates.
This supplies the comparison used when restoring the
cusp gauge.

Here is the coefficient construction and the exhaustive
linear solution. If conjugation by a word $q$ sends $n_i$
to normal coordinates $(A_qe_i,\tau_i)$, its action on
(E.13.3) has blocks
\[
 \begin{pmatrix}
 A_q&0\\
 T_q&\Lambda^2A_q
 \end{pmatrix},\qquad
 (T_q)_{(a,b),i}=2(\tau_i)_{ab}-(A_q)_{ai}(A_q)_{bi}
                 \pmod4,
 \tag{E.13.5}
\]
where $a<b$ indexes a normal commutator coordinate.
The formula follows by taking the sixth power of the
image of $n_i$ in (E.13.2) and then dividing its central
coordinate by three. Conjugation by an additional normal
word with abelian coordinate $c$ adds
$2(c_{i_1}(A_q)_{i_2i}-c_{i_2}(A_q)_{i_1i})$
to that block. These formulas compute every coefficient.

The program \texttt{de\_class2\_eight\_fibres.py} reduces
the $9508$ by $4754$ system modulo two. Its rank is $3927$
and its nullspace has dimension $827$, including the
multiplier variable. For a particular solution $x_0$ and
a binary nullspace matrix $K$, write every possible lift
as $x=x_0+Kz+2y$. The second-layer equation is exactly
\[
 A_0y+\frac{A_4K}{2}z
         =\frac{b_4-A_4x_0}{2}\pmod2,
 \tag{E.13.6}
\]
where the numerators are even and are reduced modulo four
before division. This system has coefficient rank $3953$
and is consistent in each of the 24 cases. It retains all
$827$ nullspace parameters, so it tests the complete fibre
above each fixed affine homology point. A solution with
$h=0$ is retained in every case. Direct substitution verifies
both order equations and the two commutator-image rows.

We then form actual words at the common refinement. Let
$k_2$ be the retained $\mathbb Z/4$ correction, and let
$k_3$ be the three-primary correction already retained in
E.11. Set
\[
 k=\prod_i n_i^{6(k_2)_i}
       \prod_{i<j}[n_i,n_j]^{3(k_2)_{ij}+4(k_3)_{ij}}.
\]
Its exponent-eight image is $k_2$, and its exponent-three
image is $k_3$. Since $H((k_2)_i)$ is divisible by four,
the same unimodular two-column calculation as in E.11 gives
an explicit product $j$ of twenty-fourth powers of basis
words satisfying $\operatorname{ab}_F(kj)=0$. It is trivial
in both factors of $J_{8,3}$. Thus
\[
                  f=sx^gS_0k x^{-g}j\in F'
 \tag{E.13.7}
\]
is a single actual word for the common shadow.

The final program \texttt{de\_fibre\_eight\_crt\_verify.py}
substitutes both complete braid hexagons for (E.13.7) in
both faithful collectors: there are $96$ successful
component checks. The exponent-eight collector also checks
sixteen block actions against literal expanded conjugate
words. It keeps the section comparison in the Schreier
calculation: if the short representative is $n(c)$ times
the spanning-tree representative, the latter's central
action is $\tau_{\rm short}-c\wedge A_q$. This is the
inner correction required by (E.13.2).

The new corrections preserve the verified $G_6$ homology
shadow. Both 97-dimensional Frattini ranks remain maximal,
so generation follows on the two primary kernels. The
actual commutator word, unit multiplier, complete hexagons
and generation prove the assertion. The JSON retains all
literal base words and the factored correction and adjustment
words. The quotient still kills $\gamma_3N$; the explicit
$A_5$ comparison at the end of E.12 therefore applies to
this refinement as well. No cofinality conclusion follows
from these positive finite fibres.
\end{proof}

\subsection{The exponent-sixteen fibres and the two-adic lifting argument}

\begin{statement}{E.14}{Proved}{Every CM matrix fibre survives the common exponent-sixteen refinement}
Let $J_{16,3}=(\gamma_3N\,N^{16})\cap(\gamma_3N\,N^3)$.
Its inverse image in $B_3$ is a finite-index normal subgroup
contained in $PB_3$. Every one of the 24 admissible CM matrix
pairs of E.1 occurs as a gentle shadow at this refinement,
with multiplier one modulo $576$.
\end{statement}
\begin{proof}
The exponent-sixteen collector has coordinates
$a\in(\Z/16)^{97}$ and $z\in(\Z/8)^{4656}$, with the product
in (E.13.2). The image of $N'N^6$, retaining $N^{\rm ab}/3$,
has normal coordinates
\[
 k=\prod_i n_i^{6u_i}\prod_{i<j}[n_i,n_j]^{3v_{ij}},
 \qquad (u,v)\in(\Z/8)^{4753},
\]
and exact product
\[
 (u,v)(u',v')=(u+u',\,v+v'-12(u_ju'_i)_{i<j}).
 \tag{E.14.1}
\]
The correction group is therefore nonabelian at this level;
the program retains its full product and inverse.

For each affine homology point of E.11 use the same ungauged
$S_0=\beta s\beta^{-1}$ and $R_0=S_0t$. With $\lambda=1+72h$,
the complete equations are
\[
 (S_0k)^2=1,\qquad
 (k^{-1}S_0t^{1+72h})^3=1,\qquad Hu=0\pmod8.
 \tag{E.14.2}
\]
Here $H$ is the two-row integer matrix of E.13.
The last condition is exact: the common kernel has
$F$-abelian image $48L_N$, and the correction has image
$6B_NHu$.

Write $x=(u,v,h)$ and let $\mathcal F(x)$ be the two
normal-coordinate order errors followed by $Hu$.
These are integer polynomials of degree at most two.
Their diagonal quadratic coefficients are divisible by six
and their mixed quadratic coefficients by twelve.
Indeed the sixth-power contribution of an automorphism
with abelian matrix $A$ and generator cochain $\tau$ is
\[
 2u_i\tau_i-(6u_i^2-u_i)A_{ai}A_{bi}
\]
in a central coordinate; reordering contributes multiples
of $12u_i u_j$. Cusp powers have the same divisibility,
since $12\binom h2=6h^2-6h$. Only linear homology
expressions are substituted into these formulas.
Thus $\mathcal F$ modulo four is exactly the affine system
of E.13.

Its homogeneous solutions are represented by
\[
 Vt+2Kw\pmod4,
 \qquad t\in\F_2^{801},\quad w\in\F_2^{827},
 \tag{E.14.3}
\]
where $K$ is the entire binary kernel and $V$ lifts the
kernel of the second-layer obstruction. Every column is
verified modulo four. The projection of $V\bmod2$ to
the 97 homology coordinates and the multiplier has rank
15. Select 15 columns spanning that projection and subtract
their combinations from the other columns. The remaining
786 columns have even homology and multiplier coordinates.
Twice a selected column is already represented by $2K$,
so this change preserves all solutions in (E.14.3).

Every direction with even homology and multiplier coordinates
changes $\mathcal F$ modulo eight linearly and independently
of the base point: the stated divisibilities kill its
quadratic cross terms. Hence all 786 remaining lower
parameters, all 827 upper parameters, and the arbitrary
correction $4y$ enter linearly; only the 15 selected bits
enter quadratically. The program
\texttt{de\_class2\_sixteen\_fibres.py} retains the full
linear system and evaluates all $32768$ choices of those
bits. In each of the 24 cases the linear correction rank
is 3953, and every constant, linear, and quadratic coefficient
vanishes modulo its image. All $32768$ choices survive.
A solution with $h=0$ is retained in every case and
substituted into both literal order equations and $Hu$.

For these two-primary coordinates and the three-primary
central correction $k_3$ of E.11, the common actual word is
\[
 k=\prod_i n_i^{6u_i}
       \prod_{i<j}[n_i,n_j]^{3v_{ij}+16(k_3)_{ij}}.
 \tag{E.14.4}
\]
The unimodular two-column calculation of E.11 supplies a
product $j$ of forty-eighth powers with
$\operatorname{ab}_F(kj)=0$. It is trivial in both factors.
Restoring the recorded cusp gauge gives
$f=sx^gS_0kx^{-g}j\in F'$.
The program \texttt{de\_fibre\_sixteen\_crt\_verify.py}
substitutes both full braid hexagons in both faithful
collectors, with 96 successful component checks.
Normal-block inversion retains the central term $-a_i a_j$.
Twelve additional literal conjugation checks include a
nontrivial change of section. Both Frattini ranks remain
97 because the corrections preserve the $G_6$ homology
point. This proves the full shadow and generation conditions.

The separate \texttt{de\_class2\_sixteen.py} finds no
lift of the 24 particular words retained in E.13. Its
projection fixes their upper corrections. Calculation
(E.14.3) permits all corrections over the homology point
and supplies the replacements above. These exact projection
maps explain the scopes of the two calculations.
\end{proof}

\begin{statement}{E.15}{Proved}{Every admissible CM pair survives the entire relative class-two tower}
For each of the 24 admissible CM matrix pairs there is a
compatible family of gentle shadows at the inverse images
in $B_3$ of
\[
             J_j=\gamma_3N\,N^j\qquad(6\mid j),
 \tag{E.15.1}
\]
ordered by divisibility, with multiplier $\lambda=1$.
All twelve extra matrix fibres therefore survive this
entire tower. The assertion specifies matrix fibres;
it does not assert extension of the particular normalized
words retained in E.13 or E.14.
\end{statement}
\begin{proof}
We construct the two-primary component first. Let
$\mathcal N_2$ be the free class-two pro-two group on
the 97 basis elements, in the extension of $Q$ obtained
by completing $N$ in this quotient. Let
$\mathcal U\subset\mathcal N_2$ have even abelian
coordinates. Multiplication by three is invertible over
$\Z_2$, so the elements $n_i^6,[n_i,n_j]^3$ give its
normal coordinates $(u,v)\in\Z_2^{4753}$, with product
(E.14.1). The polynomial bijection
\[
 Z=(u,z),\qquad z_{ij}=v_{ij}+6u_i u_j
\]
puts its operations in the form
\[
 Z*Z'=Z+Z'+\tfrac12[Z,Z'],\qquad
 [Z,Z']=(0,12u\wedge u').
 \tag{E.15.2}
\]
Every coefficient is integral. We use $\exp Z$ and
$\log$ for this explicit group-coordinate correspondence;
(E.15.2) justifies that notation at the prime two.

Conjugation by $S_0$ is an integral Lie automorphism $M_S$
with $M_S^2=1$. The first order equation is equivalent to
\[
                       G_S(Z)=(1+M_S)Z=0.
 \tag{E.15.3}
\]
In fact its logarithm is
$G_S(Z)+\frac12[G_S(Z),Z]$; the map
$Y\mapsto Y+\frac12[Y,Z]$ is an integral unipotent
automorphism. The target of $G_S$ is the saturated
submodule $\operatorname{Fix}(M_S)$.

Put $E=\log(R_0^3)$ and $r_*=\exp(-E/3)R_0$.
Here $R_0^3\in\mathcal U$, division by three is integral,
and $R_0$ commutes with its cube; thus $r_*^3=1$.
Write $M$ for its Lie action and
$P=(1+M+M^2)/3$ for the integral projection onto
$\operatorname{Fix}(M)$. For $k=\exp Z$, keep the
multiplier equal to one and write
\[
 R=k^{-1}S_0t=\exp(W)r_*,\qquad C=\log(R^3).
\]
The coordinates of $W$ and $C$ are integral polynomials
of degree at most two in $Z$, by (E.15.2).
The second order equation is equivalent to
\[
                              G_R(Z)=PC=0.
 \tag{E.15.4}
\]
To prove this integrally, the homology of $C$ is the
norm of the homology of $W$, so it is fixed by $M$ and
equals that of $PC$. Since $R^3$ commutes with $R$,
$(M-1)C=-[W,C]$. On $\ker P$ the inverse of $M-1$
is $-(M+2)/3$. With $Q_0=1-P$ this gives
\[
             C=PC+\frac{M+2}{3}Q_0[W,PC].
 \tag{E.15.5}
\]
The bracket uses only homology. Thus $PC=0$ is equivalent
to $C=0$, and (E.15.5) is an integral linear embedding
of the error $PC$, with left inverse $P$.

There is only one further independent commutator equation.
The integral matrix $H$ satisfies
\[
 H A_S=\begin{pmatrix}-2&-1\\3&2\end{pmatrix}H,
 \tag{E.15.6}
\]
where $A_S$ acts on $N^{\rm ab}$. Equivalently, apply
$B_N^{-1}$ to the interchange of the two coordinates
of $F^{\rm ab}$. The first row of $H(1+A_S)u$ is
$-(H_1+H_2)u$. Hence (E.15.3) and $H_2u=0$ imply
both commutator rows, and the converse holds. The exact
compressed system is the integral polynomial map
\[
 G:\Z_2^{4753}\longrightarrow
 \operatorname{Fix}(M_S)\oplus\operatorname{Fix}(M)\oplus\Z_2,
 \qquad G=(G_S,G_R,H_2u).
 \tag{E.15.7}
\]
The target lattices are primitively embedded. The
normal-coordinate changes and (E.15.5)--(E.15.6) recover
the full error vector from $G$, with primitive linear
part at a zero. At a point whose errors are divisible
by eight the same derivative comparison holds modulo
four, since the additional derivative terms multiply
those errors.

Here are the exact target and derivative ranks.
The cycle actions have traces $-7,-5,-5$ for $s,r,r^2$.
The exterior-square traces are $-24,15,15$, respectively,
from
$\operatorname{tr}(\Lambda^2A)=
((\operatorname{tr}A)^2-\operatorname{tr}(A^2))/2$.
Central cochains do not change a trace. It follows that
\[
 \operatorname{rank}\operatorname{Fix}(M_S)=2361,
 \quad \operatorname{rank}\operatorname{Fix}(M)=1591,
 \quad \operatorname{rank}\operatorname{target}(G)=3953.
 \tag{E.15.8}
\]
The program \texttt{de\_two\_hensel\_tangent.py}
checks these traces and (E.15.6). At each of the 24
points from E.14 it computes the full integral Jacobian
modulo four. The program \texttt{de\_two\_hensel\_fixed.py}
then deletes the multiplier column, recomputes the entire
binary kernel, and verifies all the ranks below for the
resulting 4753-variable system. Its binary kernel has
dimension 826. The divisibilities in the proof of E.14
give
\[
 \mathcal F(x+2e_i)-\mathcal F(x)=2J_i\pmod8
\]
and make $J\bmod4$ independent of $x$. The central
columns are the exterior-square norm matrices. All
homology columns are retained, with
five additional complete direction checks per point.
For the full binary kernel $K$, the exact ranks are
\[
 \operatorname{rank}_{\F_2}(J\bmod2)=3927,\qquad
 \operatorname{rank}_{\F_2}
 \left(J\bmod2\ \middle|\ \frac{JK}{2}\bmod2\right)=3953.
 \tag{E.15.9}
\]
The primitive comparison with (E.15.7) identifies these
with the ranks for $dG$. Its target has rank 3953, so
its elementary divisors consist of 3927 units and 26
elements of valuation one. Its image contains twice
the target lattice.

The lifting can now be given explicitly. At the integral
point $x_0$ retained in E.14, $G(x_0)$ is divisible
by eight. Choose a right inverse of $dG_{x_0}$ over
$\Q_2$ carrying the integral target into
$\frac12\Z_2^{4753}$; the elementary divisors give
such an inverse. The Newton correction belongs to
$4\Z_2^{4753}$. The new error is divisible by sixteen,
because $G$ is integral of degree at most two, and its
derivative is unchanged modulo four. Repeating, an
error of valuation $v\geq3$ is replaced by one of
valuation at least $2v-2$. The corrections converge to
an exact zero of $G$, with multiplier still equal to one.
Equations (E.15.3)--(E.15.6) give both exact orders
and zero $F$-abelianization for the correction $k$.
The retained cusp gauge commutes with $t$ and makes
the normalized $f$ have zero abelianization, by the
formula (E.15.12) below. The reduction on $N^{\rm ab}/2$ has
rank 97, so the pro-two Frattini argument gives an
automorphism at every finite level.

We give the gluing before applying the cusp gauge.
For the same CM matrix pair, take the ungauged
three-primary action of E.12, and take the identity
action at every other prime. In each component the
images of $s,r,t$ in $Q$ are $s,r,t$: the matrix
representatives $\beta,A$ centralize $s,r$ in $Q$,
and all their corrections belong to the normal kernel.
Thus these ungauged actions have the same identity
action on $Q$. The precise comparison of their target
groups is
\[
 \varprojlim_j B/(\gamma_3N\,N^j)
 \simeq
 \prod_{p,Q}\varprojlim_a B/(\gamma_3N\,N^{p^a}),
 \tag{E.15.10}
\]
where the product is the fibre product over $Q$.
On the normal kernel this is the Sylow decomposition
of the class-two completion of $N$: every finite
nilpotent quotient is the product of its Sylow groups.
For surjectivity at finitely many levels, lift their
common $Q$ coordinate and adjust independently in
these normal Sylow groups. Compactness gives (E.15.10).

We now specify one gauge in this common group. Let
$g_2\in\Z_2$ be the integral gauge retained in E.14.
For the limiting three-primary conjugator $\beta_3$,
let $g_3\in\Z_3$ be the difference of its second and
first $F$-abelian coordinates. The E.12 corrections
preserve these coordinates modulo three, so the
reduction of $g_3$ is that of the recorded integral
gauge. The image of $N$ in $F^{\rm ab}$ has basis matrix
$B_N=\left(\begin{smallmatrix}6&4\\0&2\end{smallmatrix}\right)$;
therefore every integral change in this difference is
even. The two-primary residue of the three-primary
gauge can accordingly be kept equal to that of $g_2$.
Choose $g\in\widehat\Z$ with these two primary
components and with $g_p=0$ for $p\ne2,3$.
The element $x=t^{-2}$ has order six in $Q$.
Consequently $x^g$ restricts to the specified cusp
gauge in the two- and three-primary extensions.
Apply its inner conjugation to the glued ungauged
action, obtaining $S',R'$ with $S'R'=t$.

Put $f=sS'$. Its common $Q$ coordinate is
\[
             q(f)=s x^{\bar g}s x^{-\bar g},
             \qquad \bar g=g\bmod6.
 \tag{E.15.11}
\]
In particular $f$ need not lie in the image of $N$.
The independent matrix program
\texttt{de\_adelic\_gauge.py} verifies the identity
ungauged $Q$ action and (E.15.11) for all 24 pairs.
For $\bar g=0,4,2$, respectively, its canonical matrix
representatives in $\mathrm{SL}_2(\Z/12)/\{\pm1\}$ are
\[
 I,\qquad
 \begin{pmatrix}1&4\\4&5\end{pmatrix},\qquad
 \begin{pmatrix}1&8\\8&5\end{pmatrix}.
\]
The exact abelianization formula is
\[
 \operatorname{ab}_F(sx^g Sx^{-g})
       =\operatorname{ab}_F(sS)+(-g,g).
 \tag{E.15.12}
\]
At two and three the selected gauges make this zero.
At every other prime the ungauged action is identity,
and $g_p=0$, so it is zero there also. In those
components the final action is inner by $x^g$; its
image can have nontrivial prime-to-$p$ order dividing
six, exactly as (E.15.11) requires. Hence $f$ has zero
full $F$-abelianization. Since $\gamma_3N\subset F'$,
the kernel of the resulting map to
$\widehat F^{\rm ab}=\widehat\Z^2$ is the image of
$\widehat F'$, and $f$ belongs to that image.
Both order equations hold in every component. The
exact hexagon equivalence of E.9 now supplies compatible
gentle shadows. Generation holds on each normal Sylow
kernel by its Frattini map at two and three and by
the specified inner maps at the other primes; it also
holds on $Q$. Each finite map is therefore an
automorphism. The two- and three-primary reductions
recover the prescribed CM matrix pair, proving the family.

Every subgroup in (E.15.1) contains $\gamma_3N$.
The map to $A_5$ in E.5 sends that subgroup onto
$A_5$, so its common refinement factors through none
of these quotients. This explicit comparison leaves
the cofinal lifting problem of E.2 unresolved.
\end{proof}

\begin{statement}{E.16}{Proved}{Simultaneous extension through every nonabelian simple quotient}
Let $\Sigma$ be any finite collection of epimorphisms
$\theta:B\twoheadrightarrow S_\theta$ to nonabelian finite
simple groups, and put
\[
 K_{j,\Sigma}=\gamma_3N\,N^j\cap
                \bigcap_{\theta\in\Sigma}\ker\theta,
 \qquad 6\mid j.
 \tag{E.16.1}
\]
Every one of the 24 CM matrix pairs has a compatible family
of gentle shadows at the inverse images of these subgroups
in $B_3$, with multiplier one, as $j$ increases by
divisibility and $\Sigma$ increases by inclusion.
Thus the class-two families of E.15 extend simultaneously
through all these nonabelian simple quotients.
For each fixed epimorphism $B\twoheadrightarrow R$ to
an arbitrary finite perfect group, they also extend
compatibly in $j$ through
$\gamma_3N\,N^j\cap\ker(B\to R)$.

The finite solvable quotient
$B/(\gamma_4N\,N^5)$ is not a quotient, through the natural
map from $B$, of any $B/K_{j,\Sigma}$. This gives a precise
remaining finite comparison outside this enlarged family.
\end{statement}
\begin{proof}
Write $E_j=B/(\gamma_3N\,N^j)$ and let $R_\Sigma$ be the
image of $B$ in $\prod_{\theta\in\Sigma}S_\theta$.
The latter is a product of nonabelian simple groups,
with repeated diagonal factors identified. To prove this,
adjoin the simple factors one at a time. The kernel of
projection to the previous factors, viewed in the new
simple group, is either the whole group or trivial.
In the first case the new group is a direct product.
In the second, the new factor is a simple quotient of
the preceding product and is the graph of an isomorphism
from one of its simple factors. This proves the claim
by induction. In particular $R_\Sigma$ is perfect.

The group $Q$ is solvable. Indeed,
$\mathrm{SL}_2(\Z/12)$ is the product of its reductions
modulo four and three. The kernel of reduction
$\mathrm{SL}_2(\Z/4)\to\mathrm{SL}_2(\F_2)$ is an
elementary abelian group of order eight, and the target
is $\mathfrak S_3$. The other factor has a normal
quaternion group of order eight with cyclic quotient
of order three. Passing to the quotient by $\{\pm I\}$
preserves solvability. Since the kernel of $E_j\to Q$
is nilpotent, $E_j$ is solvable as well.

The image of $B$ in $E_j\times R_\Sigma$ projects
onto both factors. For a subdirect product $H\subset U\times V$,
put $U_0=\{u:(u,1)\in H\}$ and
$V_0=\{v:(1,v)\in H\}$. Sending one coordinate of a
pair in $H$ to the other defines an isomorphism
$U/U_0\simeq V/V_0$. Here that common quotient is
both solvable and perfect, hence trivial: the derived
series of a solvable perfect group is simultaneously
constant and eventually trivial. Therefore
\[
                    B/K_{j,\Sigma}\simeq E_j\times R_\Sigma.
 \tag{E.16.2}
\]
The same argument applies to the image of $F'$.
First $F$ maps onto $R_\Sigma$, since $B/F\simeq
\mathfrak S_3$ and a quotient of a perfect group that
is solvable is trivial. Then $F'$ also maps onto
$R_\Sigma$, which is perfect. Its image in $E_j$ is
solvable, so its joint image is the product of its two
images.

For each E.15 shadow use its automorphism on $E_j$
and the identity on $R_\Sigma$. Under (E.16.2) the
conjugator is $(f_j,1)$, which is in the image of
$F'$ by the preceding paragraph. Both hexagons hold
coordinatewise with multiplier one, and the product
map is an automorphism. These constructions commute
with every transition map. Equivalently, compactness
applied to the joint images of $\widehat F'$ gives
one lift of the compatible pair $(\{f_j\},1)$.
This proves the claimed compatible family.
For one fixed arbitrary perfect target $R$, the same
subdirect-product and $F'$ arguments apply directly
to $E_j\times R$ and give the additional assertion.
For a collection of such targets the proof above
uses perfection of their \emph{joint image}; the
simple-factor induction proves exactly that property
for the simultaneous family in (E.16.1).

For the final comparison set $L=\gamma_4N\,N^5$.
It is characteristic in $N$, hence normal in $B$,
and has finite index because $N$ is finitely generated
and its class-three exponent-five quotient is finite.
The quotient $B/L$ is solvable. Every homomorphism
from (E.16.2) to a solvable group kills $R_\Sigma$,
since its image is perfect. Therefore a factorization
of the natural map $B\to B/L$ through (E.16.2)
would factor through $E_j$ and kill $\gamma_3N$.

Here is an explicit word that it does not kill.
Send the first two free generators of $N$ to
\[
 a=I+E_{12}+E_{34},\qquad b=I+E_{23}
                    \quad\hbox{in }U_4(\F_5),
\]
and send the others to the identity. The target has
class three and exponent five: commutators add
superdiagonal degrees, and $(I+A)^5=I$ for a strictly
upper triangular $4\times4$ matrix $A$ in
characteristic five. With $[a,b]=aba^{-1}b^{-1}$,
direct multiplication gives
$[[a,b],a]=I+2E_{14}\ne I$.
The program \texttt{de\_solvable\_barrier.py}
checks these literal matrices. Hence
$[[n_0,n_1],n_0]\in\gamma_3N\setminus L$,
contradicting that factorization. This proves the
claimed finite comparison and the exact scope of
the enlarged family.
\end{proof}

\subsection{A four-strand pentagon constraint on the CM fibres}

\begin{statement}{E.17}{Proved}{The complete Ferrari-homology pentagon test and its affine fibres}
Put $H_3=\F_3^2\rtimes\mathrm{SL}_2(\F_3)$, with the
standard left action on column vectors. The assignments
\[
 \begin{aligned}
 \rho_4(\sigma_1)&=(0,A),&
 A&=\begin{pmatrix}1&1\\0&1\end{pmatrix},\\
 \rho_4(\sigma_2)&=(0,B),&
 B&=\begin{pmatrix}1&0\\-1&1\end{pmatrix},\\
 \rho_4(\sigma_3)&=(e_1,A)
 \end{aligned}
 \tag{E.17.1}
\]
define an epimorphism $B_4\to H_3$.
Let $L_4=\ker\rho_4\cap PB_4$. Then $L_4$ is normal
in $B_4$, is contained in $PB_4$, and
$|B_4/L_4|=5184$.

The complete pentagon in $PB_4/L_4$ is a well-defined
condition on the $G_6$ conjugator coordinate used in E.1.
Among the 24 particular $G_6$ representatives retained
in E.14, exactly nine satisfy it. They have indices
\[
             3,4,5,7,8,14,15,17,23;
 \tag{E.17.2}
\]
the extra-coset indices among these are $4,7,15$.
The other fifteen particular representatives cannot
be projections of full $\widehat{\mathrm{GT}}$ elements.

Every one of the 24 \emph{whole CM matrix fibres}, however,
contains a gentle $G_6$ shadow with multiplier one
whose actual commutator word satisfies this pentagon.
In the 18-parameter three-primary presentation of E.10,
before quotienting by cusp conjugations, its equation
is one nonzero affine linear equation. Its solution
space has dimension 17 and contains points with full
generation rank 97 and with vanishing class-two cusp
obstructions of E.11. This is a finite pentagon test;
no lift through all four-strand quotients is asserted.
\end{statement}
\begin{proof}
We use the multiplication
$(u,M)(v,T)=(u+Mv,MT)$ in $H_3$.
The three assignments in (E.17.1) satisfy both adjacent
braid relations and the relation between the distant
generators, by direct multiplication. Their Ferrari
kernel elements
\[
 \xi_1=\sigma_3\sigma_1^{-1},\qquad
 \xi_2=\sigma_1\sigma_2\sigma_1
                  \sigma_3^{-1}\sigma_2^{-1}\sigma_1^{-1}
\]
map to $(e_1,I),(e_2,I)$, while $A,B$ generate
$\mathrm{SL}_2(\F_3)$. This proves surjectivity.
These are the actual free Ferrari generators and their
homology action: the source comparison is
\cite[(2.9) and (4.15)]{DE_NakamuraTsunogaiYasuda2010}.
The program \texttt{de\_ferrari\_pentagon.py} checks the
literal relations and enumerates the joint image with
the adjacent transpositions in $\mathfrak S_4$;
its order is 5184. This is precisely $B_4/L_4$.

Write $a_{ij}$ for the standard pure Artin generators.
We evaluate the complete group equation
\[
 \begin{aligned}
 &f(a_{12},a_{23}a_{24})f(a_{13}a_{23},a_{34})\\
 &\qquad=f(a_{23},a_{34})
       f(a_{12}a_{13},a_{24}a_{34})f(a_{12},a_{23}).
 \end{aligned}
 \tag{E.17.3}
\]
Each displayed pair defines an actual homomorphism
$F_2\to PB_4$, followed by $\rho_4$.
All five factors are pure, so the permutation component
of (E.17.3) is already the identity. We also retain
the evaluation at
$((\sigma_1\sigma_2)^3,\sigma_3^2)$, which is the
conjugator in the standard Drinfeld image of
$\sigma_3$.

The simultaneous images of the 97 literal generators
of $N$ under these six homomorphisms commute and have
exponent three. The program checks every generator
cube and all $\binom{97}{2}=4656$ commutators.
It follows that each homomorphism kills $N'N^6$.
Thus all six evaluations factor through the actual
quotient $F_2/(N'N^6)$, which is the pure free factor
in $G_6$. Corrections by the whole kernel have zero
image in these targets. Consequently a failure of
(E.17.3) at the computed representative is a failure
for every lift of that $G_6$ conjugator.
The retained E.14 word differs from the base word
evaluated here by elements of $N'N^6$, so this
comparison applies to the E.14 representatives
themselves. Complete multiplication of all five
factors gives exactly (E.17.2).

Here is the exhaustive affine reduction for whole
CM matrix fibres. Let $P=\binom{B_0}{A_0}$ be the
$194\times18$ parameter matrix retained in E.10.
A parameter $z\in\F_3^{18}$ changes the two
conjugators by $n(B_0z),n(A_0z)$. For an actual
integral change $b$ in the first conjugator, put
\[
 \ell(b)=\frac{
       \operatorname{ab}_{F,2}(n(b))-
       \operatorname{ab}_{F,1}(n(b))}{2}\pmod3.
\]
The numerator is even by the matrix $B_N$ of E.9.
The new gauge is $g+2\ell(b)$ modulo six.
If $f_\beta=s\beta s\beta^{-1}$ denotes the
ungauged base conjugator, the actual new word is
\[
 f_z=y^{g+2\ell(b)}
       (s n(b)s) f_\beta n(b)^{-1}
                  x^{-g-2\ell(b)}.
 \tag{E.17.4}
\]
This identity follows by inserting $s^2=1$ and
$sxs=y$; it is the precise gauge comparison used
in the calculation.

Embed each affine group in $3\times3$ matrices.
On the commuting exponent-three images of $N$,
the formulas
\[
 \log(I+D)=D+D^2,\qquad
 \exp(L)=I+L+2L^2\quad\text{over }\F_3
\]
are inverse and turn products into sums. Their
nilpotent matrices have cube zero, and every
commutator is zero. A commuting family of nilpotent
endomorphisms admits a common complete invariant flag:
choose a nonzero common kernel vector, pass to its
quotient, and repeat. In that basis these matrices
are strictly upper triangular, so every product of
three is zero. Expanding the two displayed polynomials
therefore proves $\log(uv)=\log(u)+\log(v)$ for
these commuting matrices, including characteristic three.
The joint linear map
recording the six logarithms of $n(b)$, those of
$s n(b)s$, and $\ell(b)$ has rank two on the
18 parameter columns. The program
\texttt{de\_ferrari\_fibres.py} retains that exact
matrix, its two-column section, and its
16-dimensional kernel. It enumerates all nine
values of the two-dimensional image. Three
independent literal-word substitutions per fibre
check (E.17.4) against the logarithmic evaluation.

For every one of the 24 fibres, the pentagon
residual on those nine values is an affine
linear map of rank one. Every value is retained
in the JSON certificate, and direct substitution
checks its equality with the retained linear
map. Exactly three values satisfy the pentagon.
Each also gives an invertible induced action
on Ferrari homology; its two columns are the
images of $\xi_1,\xi_2$ under the three
Drinfeld generator formulas. The program checks
both intertwining relations and all braid
relations after this substitution.

Finally, in the remaining 16 parameters the
program retains an actual point in every fibre
at which the five class-two cusp equations of
E.11 vanish and the complete $97\times97$
generation matrix is invertible. These are
positive witnesses, not exclusions from a
failed search. Their three-primary changes
are represented by four times the corresponding
integral $N$ words, so their two-primary
homology component is unchanged. The literal
cusp relation in $G_6$, both generation ranks,
the exact integral gauge, zero $F$-abelianization,
and all five pentagon factors are checked again.
The resulting commutator words and every
parameter are retained. This proves nonemptiness
in all 24 whole matrix fibres and the asserted
finite scope.
\end{proof}

\begin{statement}{E.18}{Proved}{The next Ferrari central layer adds no equation on these affine fibres}
Let $\mathcal H_3$ be the class-two exponent-three free group on two
generators, in coordinates $(u,v,z)\in\F_3^3$ with multiplication
\[
 (u,v,z)(u',v',z')=
 (u+u',v+v',z+z'+2(uv'-vu')).
\]
The standard action of $\mathrm{SL}_2(\F_3)$ on $(u,v)$ fixes $z$.
Put $\widetilde H_3=\mathcal H_3\rtimes\mathrm{SL}_2(\F_3)$.
The assignments
\[
 \sigma_1\longmapsto((0,0,0),A),\qquad
 \sigma_2\longmapsto((0,0,0),B),\qquad
 \sigma_3\longmapsto((1,0,2),A)
 \tag{E.18.1}
\]
define an epimorphism $\widetilde\rho_4:B_4\to\widetilde H_3$,
where $A,B$ are those in (E.17.1). The normal pure kernel
$\widetilde L_4=\ker\widetilde\rho_4\cap PB_4$ has index $15552$.
Reduction of the Heisenberg central coordinate gives the precise
quotient $\widetilde H_3\to H_3$ of E.17.

All six evaluations used in E.17 still factor through $G_6$.
On each of the 24 whole affine CM fibres, their joint parameter
image has dimension two, with exactly the same kernel as in E.17.
The complete pentagon has exactly the same zero set on this
image: three of its nine points. Every retained E.17 literal
witness satisfies the complete pentagon and all three braid
relations in $\widetilde H_3$, and its three braid images generate
that group. Thus this additional central layer preserves all
24 nonempty fibres and introduces no further affine equation.
\end{statement}
\begin{proof}
In the usual Baker--Campbell--Hausdorff coordinates of the free
class-two group, write $z_0$ for the central coordinate and
take the free generators to have $z_0=0$. The literal Ferrari
actions (4.15) of \cite{DE_NakamuraTsunogaiYasuda2010} have
central linear terms $(0,1)$ for $A$ and $(1,0)$ for $B$.
The exact change of coordinates
\[
                 z=z_0-u-2v
\]
removes both terms: for $\ell=(1,2)$, their rows are
$\ell A-\ell=(0,1)$ and $\ell B-\ell=(1,0)$.
Consequently the actual Ferrari generators become
$(e_1,2),(e_2,1)$, and their commutator is $(0,0,1)$.
This proves that (E.18.1) uses the actual Ferrari quotient.
Direct multiplication verifies the braid relations and the
four literal free-generator action formulas. These elements
generate the Heisenberg kernel and $\mathrm{SL}_2(\F_3)$,
proving surjectivity and the order $27\cdot24=648$.
The program \texttt{de\_ferrari\_class2\_model.py} also verifies
the Heisenberg multiplication and all 24 matrix actions
exhaustively, then enumerates the joint image with
$\mathfrak S_4$, of order $15552$.

For each of the six evaluation homomorphisms, the simultaneous
images of $N$ have order at most nine. The program
\texttt{de\_ferrari\_class2\_fibres.py} enumerates each such
subgroup and verifies that it is elementary abelian. In the
seven displayed coordinates $(a,b,c,d,u,v,z)$, subtracting
the identity gives additive coordinates on each of these
subgroups; every pair product is checked. All generator
cubes and all 4656 generator commutators are therefore
trivial. This proves the factorization through $G_6$.

Stack these additive coordinates for $n(B_0t)$ and
$s n(B_0t)s$, together with the gauge row $\ell(B_0t)$
of E.17. The retained exact matrix has rank two. Stacking
it with the old E.17 matrix still has rank two, so their
kernels in $\F_3^{18}$ coincide. Formula (E.17.4) remains
an exact literal-word identity in the larger group.
The calculation substitutes all nine values of the
two-dimensional image in every fibre. Its full
three-coordinate pentagon residual vanishes precisely
at the three points where the E.17 two-coordinate
residual vanishes. Three further literal substitutions
per fibre check the additive-coordinate evaluation.
Finally, each of the 24 saved E.17 words is substituted
directly; all five factors, all braid relations, and
generation of the order-648 group are checked. The
JSON retains every one of the $24\cdot9$ residuals,
and the compressed certificate retains the complete
parameter matrix, its section, and its kernel.
\end{proof}
\input{sections/E_ferrari_mod9_agent.tex}



